\documentclass[10pt,twocolumn]{article}

\usepackage[margin=0.72in]{geometry}
\usepackage{amsmath,amssymb,mathtools}
\usepackage{booktabs}
\usepackage{graphicx}
\usepackage{microtype}
\usepackage{natbib}
\usepackage{xcolor}
\usepackage[hidelinks]{hyperref}

\newcommand{\method}{Control-Guided Predictive Gradient Admission}
\newcommand{\shortmethod}{CGPA}
\newcommand{\todo}[1]{\textcolor{red}{[TODO: #1]}}
\newcommand{\R}{\mathbb{R}}
\DeclareMathOperator{\proj}{proj}

\title{Who Shapes the Encoder?\\
Control-Guided Prediction for Latent World Models}
\author{Anonymous Authors}
\date{}

\begin{document}
\maketitle

\begin{abstract}
Joint-embedding predictive world models learn compact latent dynamics by
predicting future representations rather than pixels.  Their prediction
objective, however, plays two roles at once: it trains a dynamics predictor
and determines which observable factors are allowed to shape the shared
encoder.  A small predictable nuisance can consequently dominate latent
geometry while anti-collapse statistics and even task-variable decodability
remain healthy.  We study this failure in action-conditioned visual world
models and distinguish information presence from information allocation.
We introduce \method{} (\shortmethod{}), an asymmetric optimization interface
that preserves the complete prediction gradient for the predictor while
selectively admitting it into the encoder according to an action-derived
control signal.  On a watermarked PushT benchmark, the method reverses the
representation from nuisance-dominated to content-dominated geometry while
retaining physical-state decodability and improving latent-space planning.
Experiments across PushT, RandGoal, TwoRooms, OGBench-Cube, and Reacher expose
an important boundary: an instantaneous control gradient is evidence about
decision relevance, not a complete certificate of it.  We therefore evaluate
soft residual admission, negative-conflict-only routing, and multi-component
control-gradient subspaces to separate directional guidance from generic
gradient attenuation.  Our results frame predictable-nuisance failure as an
optimization-allocation problem and clarify when action-only supervision can
and cannot protect planning-relevant representations.
\end{abstract}

\section{Introduction}

Joint-embedding predictive architectures (JEPAs) offer an attractive route to
visual world models: encode observations, predict future embeddings under
candidate actions, and plan using distances in the learned latent space.
LeWorldModel (LeWM) demonstrates that this recipe can be trained end-to-end
from pixels with a compact prediction objective and a statistical
anti-collapse regularizer \citep{maes2026leworldmodel}.  Yet preventing global
collapse does not ensure that the finite latent budget is allocated to the
factors needed by a planner.

The Obsessed Encoder study makes this distinction concrete
\citep{enigma2026obsessed}.  A tiny, low-entropy, temporally predictable visual
tag can occupy a disproportionate share of a JEPA representation.  The model
continues to satisfy its anti-collapse objective and can achieve low latent
prediction error, but its planner compares observations largely by the tag
rather than by task content.  Filtering static information is not a general
solution: a task goal can itself be static.  The relevant distinction is not
dynamic versus static, but whether changing a factor changes future decisions.

We begin from a simple observation.  The latent prediction loss has two
separable effects.  It must train the predictor to map a current latent and an
action to a future latent.  At the same time, its gradient through the target
and context encoders determines what the representation becomes.  These roles
need not receive identical optimization access.  We therefore preserve the
full predictive learning signal for the predictor, but treat its encoder-side
gradient as a proposal whose priority is informed by an action-derived control
objective.

This intervention is deliberately narrower than claiming that action
supervision identifies every useful physical variable.  Inverse dynamics can
make action-relevant information decodable without ensuring that it dominates
the planner-facing metric.  Conversely, a direction orthogonal to the current
inverse-dynamics gradient may encode already learned control information,
long-horizon consequences, or a goal variable absent from the short-horizon
proxy.  Our central question is therefore:

\begin{quote}
How can a predictive world model reduce nuisance takeover while retaining a
reliable learning path for dynamics not yet covered by its control guide?
\end{quote}

Our contributions are:
\begin{itemize}
    \item We separate \emph{predictor fitting} from \emph{encoder allocation}
    and diagnose predictable-nuisance failure at their optimization interface.
    \item We introduce an asymmetric, per-sample rule that preserves full
    predictor learning while admitting control-supported prediction gradients
    into the encoder and attenuating unsupported components.
    \item We show on watermarked PushT that planning failure can persist even
    when physical variables are linearly decodable, whereas changing their
    relative geometric prominence restores planning.
    \item We introduce falsifying ablations---negative-conflict-only routing,
    norm-matched scalar attenuation, shuffled guides, soft residual admission,
    and component-gradient span diagnostics---to determine whether gains arise
    from direction selection or merely reduced prediction-gradient magnitude.
    \item We evaluate the boundary between nuisance suppression and dynamics
    preservation across five task settings without using privileged physical
    state during training.
\end{itemize}

We do not claim to invent gradient surgery.  Our contribution is the failure
mechanism, the asymmetric JEPA encoder--predictor interface, and an empirical
account of which parts of predictive learning should be allowed to organize a
planner-facing representation.

\section{Related Work}

\paragraph{Latent predictive world models.}
LeWM combines next-embedding prediction with SIGReg and performs model-based
planning directly in the learned representation \citep{maes2026leworldmodel}.
We retain its model and planner and study a failure that remains invisible to
global anti-collapse statistics: a representation can be diverse yet allocate
its useful geometry to an easy nuisance.  The Enigma case study identifies
this predictable-feature takeover across modern JEPA systems
\citep{enigma2026obsessed}; our work asks how the optimization interface can be
changed without defining the nuisance in advance.

\paragraph{Action and state grounding.}
Inverse dynamics is a natural non-privileged signal because actions accompany
robot trajectories and force transitions to retain action-readable
information.  Recent physically grounded JEPA work combines inverse dynamics
with direct state alignment \citep{liu2026physicallygrounded}.  Physical state
provides a strong reference but is not always available and changes the
supervision setting.  We use only pixels, actions, and temporal organization
during training.  More importantly, we distinguish adding a control loss from
regulating how the prediction loss reshapes the shared encoder.

\paragraph{Gradient interaction.}
Gradient similarity has long been used to decide whether an auxiliary update
helps a main task \citep{du2018auxiliary}.  Auxiliary Task Update Decomposition
separates helpful, harmful, and locally neutral components
\citep{dery2021atud}.  PCGrad removes only negatively conflicting task-gradient
components and preserves orthogonal updates \citep{yu2020pcgrad}.  Bloop also
protects an orthogonal auxiliary component, using an exponential moving average
to stabilize stochastic projections \citep{hsieh2024bloop}.  These works make
it especially important not to equate orthogonality with irrelevance.

Our setting reverses the usual priority: the predictive objective is the
high-bandwidth auxiliary signal whose access to the encoder is regulated by a
control guide, while the prediction head still receives its unmodified
gradient.  Moreover, the observed failure is dominated by near-orthogonal
rather than negative-conflict updates.  We consequently include PCGrad-style
conflict removal as a necessary comparator and treat attenuation of the
orthogonal component as a testable modeling choice, not a theorem about its
semantic content.

\section{Method}

\subsection{Predictive learning and control guidance}

Let $o_t$ be an image, $a_t$ an action, and
$z_t=f_\theta(o_t)\in\R^d$ an encoder representation.  A predictor
$p_\phi$ estimates a future representation,
\begin{equation}
    \widehat z_{t+1}=p_\phi(z_t,a_t), \qquad
    \mathcal L_p=\|\widehat z_{t+1}-z_{t+1}\|_2^2.
\end{equation}
The base model combines $\mathcal L_p$ with a statistical regularizer that
prevents complete representation collapse.  We add a control objective
$\mathcal L_c$ constructed from inverse action prediction over one or more
horizons, with optional action-cycle and reachability terms.  It uses images,
actions, and trajectories, but no simulator state.

Simply optimizing $\mathcal L_p+\lambda_c\mathcal L_c$ does not determine which
objective organizes the shared metric.  We instead inspect their per-sample
gradients at the representation interface,
\begin{equation}
    g_p=\nabla_z\mathcal L_p, \qquad
    g_c=\nabla_z\mathcal L_c.
\end{equation}
Writing
\begin{equation}
    \alpha=\frac{g_p^\top g_c}{\|g_c\|_2^2+\epsilon},\quad
    g_{\parallel}=\alpha g_c,\quad
    g_\perp=g_p-g_{\parallel},
\end{equation}
separates the prediction proposal into components parallel and orthogonal to
the current control gradient.  Figure~\ref{fig:method-overview} summarizes the
asymmetric interface and the three routing variants used to test it.

Near-orthogonality is the geometric null expectation in high dimension, not
evidence of semantic irrelevance.  For independent isotropic unit vectors in
$\R^d$, the cosine has mean zero and variance $1/d$.  Consequently, a raw
cosine gate will usually be small even when neither objective is pathological.
Equation~\ref{eq:cosine-routing} must therefore be interpreted as both a
directional preference and a strong prediction-allocation bottleneck.  Our
ablations in Section~\ref{sec:orthogonality-challenge} are designed to separate
these two effects.

\begin{figure*}[t]
\centering
\definecolor{predblue}{HTML}{3274A1}
\definecolor{ctrlorange}{HTML}{D97818}
\definecolor{routegreen}{HTML}{2F8F46}
\definecolor{nuisred}{HTML}{B44747}
\definecolor{panelgray}{HTML}{F3F4F6}
\begin{tikzpicture}[
    x=1cm,
    y=1cm,
    font=\sffamily\footnotesize,
    >={Latex[length=2mm]},
    box/.style={draw=black!48, rounded corners=2pt, fill=white,
        minimum height=0.68cm, align=center, inner xsep=5pt},
    latent/.style={box, fill=predblue!9, draw=predblue!70},
    control/.style={box, fill=ctrlorange!11, draw=ctrlorange!85},
    router/.style={box, fill=routegreen!11, draw=routegreen!85,
        minimum width=1.55cm},
    flow/.style={->, line width=0.65pt, draw=black!65},
    predgrad/.style={->, dashed, line width=0.9pt, draw=predblue},
    ctrlgrad/.style={->, dashed, line width=0.9pt, draw=ctrlorange},
    routegrad/.style={->, line width=1.05pt, draw=routegreen},
    note/.style={align=center, font=\sffamily\scriptsize},
    paneltitle/.style={font=\sffamily\small\bfseries, anchor=west}
]

% ---------------------------------------------------------------------------
% Panel A: failure mechanism
% ---------------------------------------------------------------------------
\node[paneltitle] at (0,4.25) {(a) One objective, two optimization permissions};

\node[box, minimum width=1.15cm] (obsA) at (0.65,2.75)
    {frames\\$o_t,o_{t+1}$};
\node[note, text=nuisred] at (0.65,1.95)
    {predictable\\watermark};
\draw[nuisred, line width=1pt] (0.18,3.08) -- (0.42,3.08);

\node[box, minimum width=1.15cm] (encA) at (2.25,2.75) {encoder\\$f_\theta$};
\node[latent, minimum width=1cm] (zA) at (3.72,2.75) {$z_t$};
\node[box, minimum width=1.2cm] (predA) at (5.15,2.75)
    {predictor\\$p_\phi$};
\node[latent, minimum width=1cm] (zhatA) at (6.62,2.75)
    {$\widehat z_{t+1}$};
\node[box, fill=predblue!9, draw=predblue!70, minimum width=1.15cm]
    (lossA) at (5.9,1.35) {$\mathcal L_p$};

\draw[flow] (obsA) -- (encA);
\draw[flow] (encA) -- (zA);
\draw[flow] (zA) -- (predA);
\draw[flow] (predA) -- (zhatA);
\draw[flow] (zhatA) |- (lossA);
\draw[flow] (zA) |- (lossA);

\draw[predgrad] (lossA.west) to[bend left=13]
    node[note, below, text=predblue] {shapes representation} (encA.south);
\draw[predgrad] (lossA.north) to[bend left=15]
    node[note, right, text=predblue] {fits dynamics} (predA.south);

\node[note, text=nuisred, align=left] at (3.05,0.48)
    {Easy prediction can dominate the metric\\without global representation collapse.};

% Divider
\draw[black!18] (7.25,0.15) -- (7.25,4.2);

% ---------------------------------------------------------------------------
% Panel B: asymmetric interface
% ---------------------------------------------------------------------------
\node[paneltitle] at (7.7,4.25) {(b) Control-guided predictive gradient admission};

\node[box, minimum width=1.05cm] (encB) at (8.35,2.75)
    {encoder\\$f_\theta$};
\node[latent, minimum width=0.85cm] (zB) at (9.75,2.75) {$z$};
\node[box, minimum width=1.15cm] (predB) at (11.15,3.35)
    {predictor\\$p_\phi$};
\node[control, minimum width=1.15cm] (ctrlB) at (11.15,2.05)
    {control head\\$q_\psi$};
\node[box, fill=predblue!9, draw=predblue!70] (lpB) at (12.75,3.35)
    {$\mathcal L_p$};
\node[control] (lcB) at (12.75,2.05) {$\mathcal L_c$};
\node[router] (routerB) at (12.0,0.8) {admission\\$R(g_p,g_c)$};

\draw[flow] (encB) -- (zB);
\draw[flow] (zB) -- (predB);
\draw[flow] (zB) -- (ctrlB);
\draw[flow] (predB) -- (lpB);
\draw[flow] (ctrlB) -- (lcB);

\draw[predgrad] (lpB.south) -- node[note, right, text=predblue]
    {$g_p$} (routerB.north east);
\draw[ctrlgrad] (lcB.south) -- node[note, right, text=ctrlorange]
    {$g_c$} (routerB.north west);
\draw[routegrad] (routerB.west) -| node[note, below, text=routegreen,
    pos=0.28] {routed} (encB.south);
\draw[predgrad] (lpB.north) to[bend left=22]
    node[note, above, text=predblue] {full prediction gradient} (predB.north);

\node[note, align=left] at (9.25,0.42)
    {Pixels + actions only;\\no privileged state labels.};

% Divider
\draw[black!18] (13.75,0.15) -- (13.75,4.2);

% ---------------------------------------------------------------------------
% Panel C: routing geometry and falsifiers
% ---------------------------------------------------------------------------
\node[paneltitle] at (14.15,4.25) {(c) What happens to $g_p$?};

\coordinate (O) at (14.45,1.55);
\coordinate (P) at (16.75,3.35);
\coordinate (Q) at (15.25,2.85);
\draw[->, line width=1pt, draw=predblue] (O) -- (P)
    node[above right, text=predblue] {$g_p$};
\draw[->, line width=1pt, draw=ctrlorange] (O) -- (Q)
    node[above left, text=ctrlorange] {$g_c$};
\draw[dashed, draw=black!45] (Q) -- (P)
    node[midway, right, text=black!60] {$g_\perp$};
\draw[->, line width=1.15pt, draw=routegreen] (O) -- ($(O)!0.72!(Q)$)
    node[midway, left, text=routegreen] {$R$};

\node[note, anchor=west, align=left] at (14.15,1.02)
    {\textbf{Hard:} attenuate unsupported $g_\perp$};
\node[note, anchor=west, align=left] at (14.15,0.62)
    {\textbf{Soft:} $\beta g_p+(1-\beta)R$};
\node[note, anchor=west, align=left] at (14.15,0.22)
    {\textbf{Conflict-only:} preserve $g_\perp$};

% Background panels, drawn last on background layer.
\begin{pgfonlayer}{background}
    \node[fill=panelgray, rounded corners=3pt, fit={(0,-0.02) (7.0,4.55)},
        inner sep=0pt] {};
    \node[fill=panelgray, rounded corners=3pt, fit={(7.5,-0.02) (13.5,4.55)},
        inner sep=0pt] {};
    \node[fill=panelgray, rounded corners=3pt, fit={(14.0,-0.02) (17.65,4.55)},
        inner sep=0pt] {};
\end{pgfonlayer}

\end{tikzpicture}
\caption{\textbf{Separating predictor fitting from encoder allocation.}
(a) In a standard predictive JEPA, the same loss both fits the predictor and
shapes the encoder, allowing an easy nuisance to organize latent geometry.
(b) \shortmethod{} leaves the predictor's gradient unchanged but routes the
encoder-side prediction gradient using a non-privileged control objective.
(c) Hard, soft, and negative-conflict-only variants test whether suppressing
near-orthogonal prediction updates removes nuisance takeover or discards
planning-relevant dynamics.}
\label{fig:method-overview}
\end{figure*}


\subsection{Control-guided admission}

Our zero-floor cosine admission rule is
\begin{equation}
\label{eq:cosine-routing}
    R(g_p,g_c)=[\alpha]_+g_c+
    [\cos(g_p,g_c)]_+g_\perp.
\end{equation}
A positively aligned parallel component is retained in full; a conflicting
parallel component is removed; and the orthogonal component is scaled by
non-negative cosine similarity.  Crucially, Eq.~\ref{eq:cosine-routing} replaces
the prediction gradient only at the encoder interface.  Parameters exclusive
to $p_\phi$ receive the complete gradient of $\mathcal L_p$.

This construction has limited, local guarantees.  It satisfies
\begin{equation}
    g_c^\top R(g_p,g_c)=[g_p^\top g_c]_+\geq0
\end{equation}
and explicitly bounds the admitted orthogonal energy.  Under direct Euclidean
optimization of $z$, the routed prediction term therefore cannot worsen the
current control objective to first order.  This statement does not extend
automatically through the encoder Jacobian, Adam state, stochastic batches, or
long training horizons.

\subsection{A residual learning path}

The control loss gradient is a current learning demand, not a stable measure
of information importance.  Once a control component is predicted correctly,
its residual and gradient can become small even though the corresponding
representation remains essential.  To avoid treating lack of present evidence
as proof of irrelevance, we define soft admission
\begin{equation}
\label{eq:soft-routing}
    \widetilde g_p^{(\beta)}=
    \beta g_p+(1-\beta)R(g_p,g_c),\qquad \beta\in[0,1].
\end{equation}
Here $\beta=0$ recovers zero-floor cosine routing and $\beta=1$ recovers the raw prediction
gradient.  This residual path cannot guarantee that only physical information
is retained; it tests whether cosine-gated routing closes learning channels needed by
tasks such as Reacher.

\subsection{Falsifying comparators}

\paragraph{Negative-conflict-only routing.}
Our PCGrad-style comparator preserves every orthogonal component:
\begin{equation}
\label{eq:pcgrad-comparator}
\widetilde g_p=
\begin{cases}
g_p,&g_p^\top g_c\geq0,\\
g_p-\proj_{g_c}(g_p),&g_p^\top g_c<0.
\end{cases}
\end{equation}
If negative conflicts are rare and this comparator behaves like the baseline,
then conflict removal cannot explain the main effect.

\paragraph{Magnitude and correspondence controls.}
Norm-matched scalar attenuation preserves the direction of $g_p$ but matches
the norm produced by Eq.~\ref{eq:cosine-routing}.  A retention-matched shuffled
guide uses a different sample's control axis while matching the admitted norm.
Together they separate generic encoder learning-rate reduction from correct
sample-specific direction.

\paragraph{Component-gradient span.}
Let $g_{c,1},\ldots,g_{c,K}$ be gradients of individual inverse horizons,
cycle consistency, and reachability terms.  We measure both their individual
cosines with $g_p$ and the fraction
\begin{equation}
    \frac{\|\proj_{\mathrm{span}(g_{c,1},\ldots,g_{c,K})}g_p\|_2}
    {\|g_p\|_2}.
\end{equation}
This reveals whether an apparently orthogonal summed guide hides distinct
useful directions or destructive cancellation.  Routing with respect to this
span is a candidate extension, not part of the primary method until the
diagnostic supports it.

\section{Experiments}

\subsection{Questions}

Our evaluation is organized around five questions:
\begin{enumerate}
    \item Does an easy predictable tag dominate planning geometry despite
    healthy anti-collapse statistics?
    \item Does action supervision merely make physical information decodable,
    or does it also make that information load-bearing for planning?
    \item Are gains explained by negative-conflict removal, generic prediction
    attenuation, or sample-specific control direction?
    \item Can soft residual admission preserve Reacher dynamics without
    restoring watermark dominance on PushT?
    \item What happens when a static factor is decision-causal, as in RandGoal?
\end{enumerate}

\subsection{Tasks and nuisance construction}

We use PushT, RandGoal, TwoRooms, OGBench-Cube, and Reacher under the LeWM
training and latent-space planning framework.  The tagged condition adds a
small episode-constant visual watermark that is predictable but does not alter
the environment transition or expert decision.  Clean and tagged variants
share trajectories, optimization, architecture, and planner settings whenever
possible.

RandGoal is a deliberate counterexample to static-feature filtering.  Holding
physical state fixed while changing the goal changes the expert's future
actions; the goal is therefore decision-causal even when it remains constant
within an episode.  A random watermark changes neither action selection nor
transition dynamics.  A successful allocation method should suppress the
latter without erasing the former.

\subsection{Metrics and protocol}

Planning is evaluated with the same CEM configuration, candidate budget,
evaluation seeds, and checkpoint rule within each matched comparison.  We
report success rate over complete learning curves, a predeclared late-window
mean, and final performance; single 50-episode peaks are diagnostic only.
Selected comparisons use at least three training seeds and 200--500 evaluation
episodes per selected checkpoint.

Representation evaluation includes same-content and same-tag cosine
similarity, frozen probes for task variables and tag identity, prediction and
control losses, gradient cosine, negative-conflict fraction, and admitted
prediction-gradient norm.  The content--tag margin is
\begin{equation}
    m=\mathrm{cos}_{\mathrm{same\ content}}-
    \mathrm{cos}_{\mathrm{same\ tag}}.
\end{equation}
Positive $m$ indicates that task content, rather than watermark identity,
organizes the representation.  Probe accuracy is interpreted only as
decodability; it does not establish causal use by the planner.

\subsection{Matched methods}

We compare LeWM, identical-head IDM/control-only training, summed prediction
and control losses, zero-floor cosine \shortmethod{}, soft admission with
$\beta\in\{0.25,0.5,0.75\}$, the conflict-only comparator in
Eq.~\ref{eq:pcgrad-comparator}, norm-matched scalar attenuation, and a
retention-matched shuffled guide.  State alignment is discussed as a
privileged-supervision reference but is not included as a like-for-like
training baseline.

\subsection{The orthogonality challenge}
\label{sec:orthogonality-challenge}

The central methodological objection is that near-orthogonality is common in
high dimension.  A small value of $g_p^\top g_c$ therefore does not by itself
show that the corresponding predictive update is task-irrelevant.  It can
instead mean that the control guide is low-rank, that a useful control feature
is already decoded with small residual error, that different control
components cancel in their sum, or that prediction supplies dynamics needed
only over a longer planning horizon.  We consequently do not treat
orthogonality as a semantic label.  We test what operation on the orthogonal
component is actually responsible for the planning result.

Our identification strategy uses five matched interventions.  First,
\emph{parallel-only hard drop} removes the entire orthogonal component; if it
underperforms cosine admission, orthogonal predictive information is useful.
Second, the conflict-only comparator in
Eq.~\ref{eq:pcgrad-comparator} removes only genuinely negative interference; if
it behaves like LeWM while cosine admission improves tagged planning, rare
negative conflicts are insufficient to explain the gain.  Third,
norm-matched scalar attenuation tests whether the result is merely a lower
encoder-side prediction learning rate.  Fourth, soft admission sweeps the
orthogonal residual floor and traces the trade-off between nuisance resistance
and preservation of unknown dynamics.  Fifth, the component-span diagnostic
tests whether the summed control gradient appears orthogonal because distinct
horizon, inverse, cycle, and reachability gradients cancel.

These interventions yield a falsifiable interpretation.  Evidence for
directional control guidance requires cosine admission to outperform both
norm-matched attenuation and a retention-matched shuffled guide.  Evidence
that orthogonal information is sometimes necessary requires hard drop to hurt
a clean dynamics task and soft admission to repair it.  Evidence that the
watermark effect is not standard gradient-conflict resolution requires the
conflict-only rule to fail to reproduce the tagged-task gain.  We report raw
cosine and admitted-norm distributions alongside success rates so that a
method cannot appear successful merely by silently suppressing most encoder
learning.

\subsection{Predeclared interpretation}

Soft admission is considered successful only if one shared setting improves
Reacher while retaining the PushT watermark benefit; choosing a different
$\beta$ per task is not evidence for a general method.  Conflict-only routing
tests whether negative interference suffices.  The component-span diagnostic
tests the high-dimensional orthogonality objection.  Multi-seed experiments
are promoted only after these seed-zero mechanism checks pass.

\section{Results and Current Evidence}

\subsection{Watermarked PushT}

The strongest completed result is a sustained geometric reversal in the
watermarked PushT setting.  At the end of the full seed-zero run, zero-floor cosine
\shortmethod{} yields same-content cosine $0.807$ and same-tag cosine $0.176$,
for a margin of $+0.631$.  Its final planning success is $0.80$, the mean of the
last ten evaluations is $0.774$, and the best 50-episode checkpoint reaches
$0.88$.  The peak is not used as the primary estimate.

Frozen probes show why this is not simply information deletion.  Relative to
a prediction-weight control, the projection-space block-angle probe remains
essentially unchanged ($0.800$ versus $0.795$), while tag RGB $R^2$ decreases
from $0.889$ to $0.502$.  The pusher and block positions also remain strongly
decodable.  The intervention therefore changes relative geometric prominence
more clearly than raw task-variable availability.

At convergence, the mean prediction/control cosine is approximately $0.072$
and only about $0.106$ of the raw prediction-gradient norm is admitted.  This
explains both the watermark resistance and the central concern of this paper:
cosine admission may act as an aggressive encoder-side prediction bottleneck.

\subsection{Magnitude versus direction}

In matched 10k screening runs, norm-matched scalar attenuation reaches success
$0.42$, cosine routing $0.36$, and a retention-matched shuffled guide $0.28$.
The scalar result shows that a substantial fraction of the early benefit can
come from limiting prediction-gradient magnitude.  The shuffled-guide gap is
preliminary evidence that correct sample-specific correspondence also matters.
Full schedules and multiple seeds are required before assigning either effect
priority.

\subsection{Cross-task snapshot}

Table~\ref{tab:cross-task} records completed seed-zero checkpoints.  The rows
mix 10k screening and full schedules and therefore document current evidence,
not a final leaderboard.

\begin{table}[t]
\centering
\small
\caption{Current matched seed-zero planning success. ``Full'' and ``10k'' are
not compared across rows. Reacher uses the recovered historical evaluation
protocol at epoch 10.}
\label{tab:cross-task}
\begin{tabular}{llccc}
\toprule
Task & Condition & Schedule & LeWM & Cosine \shortmethod{} \\
\midrule
TwoRooms & clean  & 10k  & 0.88 & 0.90 \\
TwoRooms & tagged & 10k  & 0.60 & 0.98 \\
Cube     & clean  & full & 0.70 & 0.76 \\
Cube     & tagged & full & 0.74 & 0.84 \\
Reacher  & clean  & epoch 10 & 0.90 & 0.86 \\
\bottomrule
\end{tabular}
\end{table}

TwoRooms and Cube do not reproduce the same catastrophic tagged-baseline
collapse as PushT.  They remain useful as robustness checks, but cannot by
themselves support a universal nuisance-removal claim.  Reacher exposes the
opposite failure boundary: aggressive cosine routing can suppress predictive learning that
the planner still needs.

\subsection{Reacher and residual admission}

A one-epoch diagnostic under the historical evaluator gives success $0.78$
for LeWM, $0.70$ for cosine routing, and $0.82$ for both control-only and
$\beta=0.75$ soft admission.  Each number is based on only 50 evaluation
episodes and is therefore a screening signal, not a final result.  The full
$\beta=0.75$ Reacher run and the paired tagged-PushT guardrail are in progress.
Their joint outcome decides whether a residual prediction path repairs the
failure or simply reintroduces the predictable shortcut.

\subsection{Static but decision-causal goals}

RandGoal shows that within-episode inverse dynamics is an incomplete guide for
static task variables.  A goal can change the expert action distribution while
remaining absent from short local transition residuals.  Current action-only
routing does not yet establish a reliable improvement in this setting.  We
retain RandGoal as a boundary condition rather than hiding it: suppressing all
static information would solve the watermark benchmark for the wrong reason.

\subsection{Pending mechanism tests}

The conflict-only PushT run tests whether the relatively rare negative
gradient conflicts explain the gain.  The component-gradient diagnostic
measures individual horizons and their joint span on tagged PushT and clean
Reacher.  We will replace this paragraph with complete curves and confidence
intervals after the predeclared runs finish.  \todo{Insert conflict-only,
soft75 guardrail, and component-span results; then promote selected variants
to seeds 1--2.}

\section{Limitations and Discussion}

The control guide is neither a semantic oracle nor a complete characterization
of planning relevance.  It can exploit a shortcut, omit a static goal, or
produce a small gradient after its decoder has learned an important feature.
The hard rule consequently provides optimization priority, not a guarantee of
physical sufficiency.  Soft residual admission makes this uncertainty explicit
but may allow a nuisance to accumulate again over long training.

Our geometric and frozen-probe analyses establish representation structure and
decodability, not causal use.  A future-latent intervention that permutes
candidate action--outcome pairings inside CEM is needed to show that improved
predictions actually determine selected actions.  Likewise, current principal
results begin with seed zero and 50-episode checkpoint evaluations.  Final
claims require multiple training seeds, more evaluation episodes, and
predeclared checkpoint aggregation.

Finally, the observed watermark collapse is task dependent: tagged TwoRooms
and Cube baselines degrade less severely than tagged PushT.  This variation is
scientifically useful because it prevents the method from being presented as a
universal filter.  The broader claim is that predictive objectives can
misallocate a planner-facing representation, and that their encoder-side
optimization authority should be measured and controlled.

\section{Conclusion}

Predictive accuracy and representation utility are not interchangeable.  A
JEPA predictor may learn a visually predictable factor that is irrelevant to
future decisions, while a short-horizon control loss may fail to cover
information essential for long-horizon planning.  By separating predictor
learning from encoder allocation, \shortmethod{} provides a practical probe
and intervention at this boundary.  The resulting evidence supports a more
careful principle: use control signals to prioritize predictive updates, but
do not interpret instantaneous orthogonality as proof that an update is
useless.


\bibliographystyle{plainnat}
\bibliography{references}

\end{document}
